% Mathematical completion of the term supplier strata.
% Direct C4 derivation-list interpretation; no upstream verifyValue claim.
% Pinned source identities are recorded in source_manifest.json.
\section{Finite-demand completion of the term supplier strata}
\label{sec:term-leaf-completion}

This section addresses arbitrary arithmetic terms in the supported vector
realization, guarded value-totality, and guarded value-coherence suppliers.
It also supplies the repeated-environment domain leaf used to discharge
the final domain guard.  The target proof objects here are explicit
finite derivation lists in $\mathcal C_4$, interpreted with its primitive
recursive proof predicate.  A statement below about a translated formula
means derivability in that calculus.  It is not a claim about the
upstream \file{verifyValue} proof-table implementation.

The lower interfaces used are the explicit source-translation, requested
label, logical proof-list, and fixed relational-template constructors
listed at the end.  Once these lower interfaces have their stated
primitive recursive implementations and quantifier-free correctness
proofs, the present argument completes the term strata.  None of the
term constructors calls back into arithmetic formula realization, PA
proof transport, or the PA induction supplier.  The appearances of full
PA induction in the inspected Lean files can therefore be replaced by
the finite-demand constructions given here; that fact is established
from their actual child calls, not from their theorem names.

\subsection{Notation and the exact padding combinators}

All source formula and term operations in this section are operations on
their numerical codes.  Write $[A\Rightarrow B]$ for
\file{imp language A B}, and $A\mathbin{\&}B$ for \file{qqAnd A B}.
For requested labels $\lambda=(n,\ell,z)$ put
\[
 P_d(\lambda)=
 (\operatorname{padEnvironment}(n,d),\ell,z+\ell+d).
\]
The actual combinators in
\file{NumericPAeqValueCertificateComposition.lean} and
\file{NumericPAeqValueCertificatePair.lean} use
\begin{align*}
 c(A,B,C)&=\max\bigl([A\Rightarrow B],
               \max([B\Rightarrow C],[A\Rightarrow C])\bigr),\\
 p(A,B,C)&=\max\bigl([A\Rightarrow B],
               \max([A\Rightarrow C],[A\Rightarrow B\mathbin{\&}C])\bigr).
\end{align*}
A composition call with source antecedent $A$, intermediate formula $B$,
and conclusion $C$ asks both children at $P_{c(A,B,C)}(\lambda)$.
A pairing call with antecedent $A$ and conclusions $B,C$ asks both
children at $P_{p(A,B,C)}(\lambda)$.  After the native logical
composition it reads the requested source conclusion back at $\lambda$.
We denote these two proof-building operations by
\[
 \operatorname{Comp}(A,B,C;f,g),\qquad
 \operatorname{Pair}(A,B,C;f,g).
 \tag{T1}
\]
This notation includes their displayed label changes.  Consequently a
nested expression in these operations is also an exact specification of
all child query labels.  No certificate value occurs in any padding
argument.

Let $T_\lambda(F)$ denote source translation and
$\operatorname{Deriv}(e,T_\lambda(F))$ the direct $\mathcal C_4$
proof-list checker.  The local correctness theorem for (T1) says that
checked child proof lists at their displayed padded labels produce a
checked parent proof list at the original labels.  Coverage of each
child follows from its displayed code bound by the relevant maximum,
its syntactic validity, and the padded-environment lemma.  This is the
precise coverage obligation used below.

\subsection{An explicit finite-demand construction}

Here is the bounded form of recursion used in this section.  A query
$q$ contains a source rank $r(q)$, a fixed tuple of environment codes,
and requested labels.  Its child-query list is computed by primitive
recursive functions of $q$, independently of all child certificate
values.  Each same-family child has rank less than $r(q)$, and there
are at most two such children.  A leaf or parent proof-list constructor
is primitive recursive in the query and its supplied child lists.

For a root of rank $r$, generate the binary demand tree to depth $r+1$,
using inactive placeholders for nonexistent children.  It has at most
$2^{r+2}-1$ slots.  Each slot's query is calculated from its ancestors;
this is bounded primitive recursion.  Along an active path its rank
decreases at every edge.  Quantifier-free induction proves that no
active nonleaf remains at the depth bound.  Evaluate the resulting tree
from its last slot towards its root.  This is again bounded primitive
recursion, so it defines an actual primitive recursive proof-list
function, not an unbounded search for witnesses.

There are two quantifier-free correctness invariants.  The forward
invariant is the bounded conjunction that every generated active child
of a valid query satisfies its local validity and coverage tests.  The
reverse invariant is the bounded conjunction that every already
evaluated active slot whose input tests hold contains a checked proof
of its specified conclusion.  Both conjunctions have primitive
recursive Boolean characteristic functions.  The forward induction
uses the syntax and padding facts.  The reverse induction uses the
local proof-list constructor equation.  Thus the root correctness
theorem is proved in $\mathsf{iPRA}$, with all root parameters free.

This argument requires the actual finite child list and an actual local
constructor.  The lists and constructors for the term suppliers are
given below.  A lower supplier stratum already equipped with a total
primitive recursive proof-list function is simply an ordinary function
call in a higher local constructor.  Its input rank may be arbitrarily
large.  This distinction is essential: the bound-variable vector leaf
calls coherence on a replacement term that need not be smaller than
the original bound-variable code.

\subsection{Removing the choice in arithmetic value syntax}

The upstream \file{NumericPAeqTerm.valueSyntax} is defined by
classical choice from a finite witness-table relation.  Its inspected
cells give the following explicit syntax computation.  Write
$\operatorname{val}(p,t)$ for its primitive recursive replacement, and write
$p=(b,f,z)$ using the actual three-entry parameter packet.  Put
\[
 C(p)=(\operatorname{shiftTerms}b,
       \operatorname{shiftTerms}f,\operatorname{bvar}0),
\]
\[
 B_i(p)=(\operatorname{shiftTerms}^2b,
         \operatorname{shiftTerms}^2f,\operatorname{bvar}i).
\]
The defining cases are
\begin{align*}
 \operatorname{val}(p,\operatorname{bvar}i)&=\operatorname{variableValue}(z,b_i),\\
 \operatorname{val}(p,\operatorname{fvar}i)&=\operatorname{variableValue}(z,f_i),\\
 \operatorname{val}(p,0_{\rm term})&=\operatorname{zeroValue}(z),\\
 \operatorname{val}(p,\operatorname{succ}u)&=
   \operatorname{successorValue}(z,\operatorname{val}(C(p),u)),\\
 \operatorname{val}(p,\operatorname{add}(u,v))&=
   \operatorname{additionValue}(z,\operatorname{val}(B_1(p),u),\operatorname{val}(B_0(p),v)),\\
 \operatorname{val}(p,\operatorname{mul}(u,v))&=
   \operatorname{multiplicationValue}(z,\operatorname{val}(B_1(p),u),\operatorname{val}(B_0(p),v)).
 \tag{T2}
\end{align*}
On an invalid source term return zero, as in the upstream definition.
Here $0_{\rm term}$ abbreviates \file{qqFunc 0 0 0}, not a new
assumption about its numerical code.  The exact nonleaf queries are
$(C(p),u)$ and $(B_1(p),u),(B_0(p),v)$.  The source proves the strict
code inequalities
\[
 u<\operatorname{succ}u,\qquad
 u,v<\operatorname{add}(u,v),\operatorname{mul}(u,v).
\]
They follow directly from the tuple and term-constructor code
inequalities, so they have quantifier-free arithmetic proofs.
The construction above applied to (T2) is primitive recursive even
though the parameter packet grows along a path: the growing packets
are determined before any child output is known, and there is a
primitive recursive finite bound on the number of demand slots.

If an exact upstream table witness is wanted, store with each output
the finite table obtained by union of the child tables and insertion of
the parent record $(p,t,\operatorname{val}(p,t))$.  These are primitive recursive
finite-set operations.  The cell equations in
\file{NumericPAeqTermTableNodes.lean} and the empty, union, and insert
equations in \file{NumericPAeqTermTables.lean} prove, by the same
bounded Boolean reverse invariant, that this table is valid and
contains the parent record.  This gives explicit witnesses rather than
using choice.

The table check here really is a Boolean finite check, despite the
existential presentation of the upstream cell predicate.  Decode each
parent term into its constructor.  In a nonleaf case search only the
records $r<C$ whose membership bits are set, compare their parameter
and term fields with the required child query, and obtain their output
from the record's output field.  A binary cell uses two such bounded
record searches.  This bounds every auxiliary output witness by a
supplied finite table record; it introduces no unbounded existential
test.  Validity is the bounded conjunction of these row checks.

Uniqueness is also bounded.  For two given valid tables, take the
bounded conjunction over pairs of their records with equal parameter
and term fields, asserting equal outputs for source rank below the
current rank.  The cell equations and injectivity and disjointness of
the term constructors prove the next case; the child records belong
to those same finite tables and have lower source rank.  All record
indices are bounded by the supplied table codes.  Quantifier-free
induction therefore proves uniqueness.  The function (T2) agrees with
the upstream uniquely characterized value syntax wherever that model
function is used.  This part needs neither PA induction nor an
unbounded witness-selection principle.

\subsection{The syntax validity and substitution joins}

The supported environments supply three Boolean facts: every displayed
target term is well formed at its stated arity; the source term is well
formed at its source arity; and each free source variable index occurring
in the term is below the free-environment length.  Each fact is checked
by recursion on the finite term code or by bounded conjunction over a
finite vector.  The source-child validity implications are literal
constructor cases: a successor has one valid child, an addition or
multiplication has two, a bound variable satisfies its bound, and a
free variable satisfies its support bound.  The target value
constructors increase the child target arity by one for successor and
two for a binary operation.  The shifts in $C(p)$ and $B_i(p)$ have
exactly those increases.  Substituting the well-formed child formulas
in these fixed constructors proves value-syntax validity by the
finite-demand Boolean invariant.

For clarity, the nontrivial target-substitution join has its own exact
query list.  Let $S_w(p)$ replace every target term in the parameter
packet by its simultaneous substitution with $w$.  Let $q(w)$ be the
usual lifted substitution vector under one new binder.  The join is
\[
 \operatorname{subst}_w(\operatorname{val}(p,t))=\operatorname{val}(S_w(p),t).                    \tag{T3}
\]
For a successor, its one child query has fields
\[
 (u,s,n+1,m+1,q(w),C(p));
\]
for a binary term, the two child queries have fields
\[
 (u,s,n+2,m+2,q(q(w)),B_1(p)),\quad
 (v,s,n+2,m+2,q(q(w)),B_0(p)).
\]
The exact parameter equalities required to join the children are
\[
 S_{q(w)}(C(p))=C(S_w(p)),\qquad
 S_{q(q(w))}(B_i(p))=B_i(S_w(p))\quad(i<2).                   \tag{T4}
\]
They follow entry by entry from simultaneous substitution of a shifted
term and from the first two entries of a lifted substitution vector.
These are the equations in
\file{NumericPAeqTermTargetSubstParameters.lean}.  The leaf cases of
(T3) are substitution into the fixed variable and zero formulas;
the successor and binary cases use (T4) and substitution through their
one or two binders.  Thus (T3), with its explicit child queries, is
another instance of the bounded induction above.  The familiar
substitution/shift identity used in (T4) itself follows by the same
term-code recursion: it is immediate for a bound or free variable,
and is componentwise through each function constructor.

For arity-zero target environments, their shifts and substitutions are
the environments themselves.  The special case $w=[x]$ of (T3) is
therefore
\[
 \operatorname{substs1}_x(\operatorname{val}((b,f,\operatorname{bvar}0),t))
       =\operatorname{val}((b,f,x),t).                                       \tag{T5}
\]
This is precisely the bound-instance equality needed by the vector
bound-variable leaf.  Source substitution into an arithmetic term
has the constructor equations
\[
 \operatorname{tsub}_w(\operatorname{bvar}i)=w_i,\quad
 \operatorname{tsub}_w(\operatorname{fvar}i)=\operatorname{fvar}i,
\]
and commutes with zero, successor, addition, and multiplication.
Its support and validity are proved by the same finite recursion,
using the supplied support tests for the $w_i$.  Applying (T3)--(T5)
to the already computed substituted term requires no bound comparing
its code to the original term.  It is a call to the completed lower
syntax stratum.

Every value formula in (T2) begins with the explicit domain conjunct
for its result variable.  Its remaining payload is obtained by the
outer conjunction projection.  Thus the payload extractions in the
successor and binary coherence templates are primitive recursive;
they require neither a search for a formula nor a new existence
theorem.  The successor and binary comparison-shape equalities follow
by (T2), (T5), and the displayed source-substitution constructor
equations.  These are the exact joins used by the files
\file{NumericPAeqSupportedValueCoherenceShape.lean} and
\file{NumericPAeqSupportedVectorRealizationShape.lean}.

\subsection{Three finite guard projections}

Domain guards, isomorphism guards, and realization guards all have the
same left-associated syntax
\[
 G_0=\top,\qquad G_{k+1}=G_k\mathbin{\&}E_k.                  \tag{T6}
\]
Their entries are respectively $D(v_k)$,
$\operatorname{Iso}(v_k,w_k)$, and
\file{entry b f w a k}.  A projection query is
$(k,i,\text{environments},\lambda)$ with $i<k$; the conclusion is the
native implication $T_\lambda(G_k)\to T_\lambda(E_i)$.
If $i=k-1$, use the fixed right-conjunction projection.
Otherwise ask exactly the one child
\[
 (k-1,i,\text{same environments},\lambda)
\]
and compose its result after the fixed left-conjunction projection.
There is no label change in this native implication construction.
The rank is $k$, and the base $k=0$ has no valid query.  Translation
of (T6) is the corresponding numerical conjunction by the source
translation equation.  These two projection equations prove the
constructor correctness directly.  A quantifier-free induction on
$k$, or the bounded demand chain, supplies a primitive recursive
projection function for all three guards.

The input-domain and coherence-input projections merely prefix a
fixed number of left/right conjunction projections to one of these
guard projections, then apply the fixed conversion from native
implication to translated source implication.  Vector-input entry
projection uses the two composition calls present in
\file{NumericPAeqSupportedVectorRealizationProjection.lean}, followed
by the realization-guard projection.  These calls use (T1), so their
requested labels are explicit as well.

\subsection{Diagonal guard and repeated-environment leaves}

Let $D_k$ be the domain guard of an environment $v$ and $J_k$ its
diagonal isomorphism guard.  A diagonal query is
$(k,v,\lambda)$ with goal $[D_k\Rightarrow J_k]$.  At $k=0$ it is
the identity on $\top$.  For $k+1$, put
$x=v_k$, $d=D(x)$, $j=\operatorname{Iso}(x,x)$, and
$A=D_k\mathbin{\&}d$.  Its exact construction is
\[
 \operatorname{Pair}(A,J_k,j;
   \operatorname{Comp}(A,D_k,J_k;\operatorname{proj}_L,
                                    \operatorname{Diag}(k,v)),
   \operatorname{Comp}(A,d,j;\operatorname{proj}_R,
                                    \operatorname{IsoRefl}(x))).
 \tag{T7}
\]
Consequently the only recursive diagonal query has labels
\[
 P_{c(A,D_k,J_k)}(P_{p(A,J_k,j)}(\lambda)),
\]
and count $k$.  The other branch is the fixed parameterized
isomorphism-reflexivity template.  Child coverage is exactly the
maximum-code coverage of (T1).  This closes
\file{diagonal_guard_supply} by bounded induction on its actual
count.  No term coherence theorem is needed for (T7).

For the repeated environment $\operatorname{repeatVec}(x,k)$, its
length is $k$ and every entry below that length is $x$.  These are
the equations of the finite repetition recursion.  Hence its domain
guard is $G_0=\top$, $G_{j+1}=G_j\mathbin{\&}D(x)$.
Let $a=T_\lambda(D(x))$.  Define explicitly
\[
 E_0:\ a\to0,\qquad
 E_{j+1}=\operatorname{nativePair}(E_j,\operatorname{Id}_a).
 \tag{T8}
\]
The base proof weakens the fixed truth proof; the step pairs two
implications with common antecedent $a$.  Induction on $j$ proves
that $E_j$ derives $a\to T_\lambda(G_j)$.  Labels remain unchanged.
The native/source implication conversion then gives exactly
\file{NumericPAeqRepeatedEnvironment.domain_requested_certificate}.
This is a primitive recursive function of $x,k,\lambda$ with
quantifier-free correctness.  It does not appeal to a choice principle
for the more general premise $\forall i<k\,\exists e_i$ of the
upstream collection theorem: in this instance every $e_i$ is the
explicit same identity proof.

\subsection{Guarded value totality: complete query list}

Abbreviate
\[
 A=D(s,b,f)=\operatorname{inputDomain}(s,b,f),\qquad
 B_t=\operatorname{valueTotalityCode}(b,f,t).
\]
A totality query $(t;s,b,f;\lambda)$ asks for $[A\Rightarrow B_t]$.
For a successor $t=\operatorname{succ}u$, the construction is
\[
 \operatorname{Comp}(A,B_u,B_t;
          \operatorname{Tot}(u;s,b,f),\operatorname{SuccTot}(u;b,f)).
 \tag{T9}
\]
Thus its sole recursive totality query has unchanged environments and
labels $P_{c(A,B_u,B_t)}(\lambda)$.  For addition or multiplication
$t=\operatorname{op}(u,v)$, put $B=B_u\mathbin{\&}B_v$.  The exact
construction is
\[
 \operatorname{Comp}(A,B,B_t;
    \operatorname{Pair}(A,B_u,B_v;
                     \operatorname{Tot}(u),\operatorname{Tot}(v)),
    \operatorname{BinTot}_{\operatorname{op}}(u,v;b,f)).          \tag{T10}
\]
Both recursive queries use the labels
\[
 P_{p(A,B_u,B_v)}(P_{c(A,B,B_t)}(\lambda)),
\]
and unchanged $s,b,f$.  Their source ranks are respectively $u,v<t$.

The leaf and local constructors in (T9)--(T10) are explicit.  A bound
or free variable first projects the appropriate $D(b_i)$ or $D(f_i)$
from $A$, then composes the variable-totality template at that target
term.  For zero, pad by
$\max(B_0,[A\Rightarrow B_0])$, take the fixed inhabited-zero proof,
weaken its conclusion by $A$, convert native implication back to source
implication, and read back the labels.  The successor local constructor
first duplicates the already present domain conjunct in the child
existential value formula, then uses the fixed successor-totality
template.  The binary local constructor duplicates the domain conjunct
in each child existential, then uses the fixed addition-totality or
multiplication-totality template.  The exact additional composition
padding in both cases is (T1) with the displayed intermediate
duplicate-domain formula; there are no further totality queries.

The fixed constructors prove that the child existential output(s)
imply the parent existential output.  This is all the local inductive
step requires.  The source child support and environment-validity
facts are the constructor facts proved above.  The finite-demand
construction therefore produces the requested totality proof for
every supported arithmetic term.

\subsection{Guarded coherence and uniqueness}

Put
\[
 A=\operatorname{coherenceInput}(s,b,c,f,g),\qquad
 B_t=\operatorname{comparisonCode}(b,c,f,g,t).
\]
The coherence query $(t;s,b,c,f,g;\lambda)$ asks for
$[A\Rightarrow B_t]$.  Its successor and binary constructions are
exactly (T9)--(T10), replacing \file{Tot} by \file{Coh},
\file{SuccTot} by successor comparison, and \file{BinTot} by
addition or multiplication comparison.  Thus the same explicit label
maps apply.  The environment tuple is unchanged, and each recursive
term code strictly decreases.

For a bound variable, project
$\operatorname{Iso}(b_i,c_i)$ from $A$ and compose the variable
coherence template.  For a free variable use
$\operatorname{Iso}(f_i,g_i)$ instead; its index is below
$\min(|f|,|g|)$ by the two support tests.  For zero use the fixed
zero-comparison proof with padding and weakening as above.  For a
successor or a binary operation extract the domain payloads of the
child value predicates, apply the corresponding fixed coherence
template, and use the exact syntax joins (T2)--(T5).

These fixed requested templates themselves have no recursive arithmetic
calls.  Variable coherence uses the closed isomorphism transport
theorem; successor coherence uses closed successor-isomorphism
compatibility; the two binary cases use closed addition and
multiplication coherence.  They are instances of the lower
singleton-requested-template constructor.

Uniqueness is a separate, nonrecursive supplier.  For its query
$(t;s,b,f;\lambda)$ set
\[
 D=\operatorname{inputDomain}(s,b,f),\quad
 A=\operatorname{coherenceInput}(s,b,b,f,f),\quad
 B=\operatorname{comparisonCode}(b,b,f,f,t).
\]
Its construction is
\[
 \operatorname{Comp}(D,A,B;
              \operatorname{DiagonalInput}(s,b,f),
              \operatorname{Coh}(t;s,b,b,f,f)).                \tag{T11}
\]
The coherence query uses $P_{c(D,A,B)}(\lambda)$.
To specify the first child completely, let $J_b$ and $J_f$ be the
diagonal isomorphism guards of $b$ and $f$, and $D_b,D_f$ their domain
guards.  Define these requested-label proof providers using (T1):
\[
\begin{aligned}
 L_b&=\operatorname{Comp}(D,D_b,J_b;\operatorname{proj}_L,
                                      \operatorname{Diag}(|b|,b)),\\
 L_f&=\operatorname{Comp}(D,D_f,J_f;\operatorname{proj}_R,
                                      \operatorname{Diag}(|f|,f)),\\
 K&=\operatorname{Pair}(D,J_b,J_f;L_b,L_f),\\
 K'&=\operatorname{Pair}(D,D,J_b\mathbin{\&}J_f;
                                      \operatorname{Id},K),\\
 \operatorname{DiagonalInput}
   &=\operatorname{Pair}(D,D,D\mathbin{\&}(J_b\mathbin{\&}J_f);
                                      \operatorname{Id},K').
\end{aligned}                                                   \tag{T12}
\]
Here $|b|=s$ by environment validity.  Expression (T12), together
with (T1), specifies the full nested label maps to its two diagonal
queries.  It is the actual three-pair/two-composition proof in
\file{NumericPAeqSupportedValueUniquenessInput.lean}.
Thus uniqueness has a primitive recursive proof constructor after
coherence and the finite diagonal guards have been completed.

\subsection{Vector realization: complete query list}

For a vector query $(t;s,r,b,f,a,w;\lambda)$ put
\[
 A=\operatorname{vectorInput}(s,r,b,f,a,w),\qquad
 B_t=\operatorname{vectorComparison}(b,f,a,w,t).
\]
Its goal is $[A\Rightarrow B_t]$.  The successor and binary
constructors again have exactly the forms (T9)--(T10), now with
recursive \file{Vec} children and the fixed successor or binary
comparison template between the substituted and original value
predicates.  The only same-family child queries replace $t$ by its
immediate source child or children, retain all six environment fields,
and use the same displayed composition/pair padding.  Their ranks
strictly decrease.

The free-variable leaf is
\[
 \operatorname{Comp}(A,D(r,b,f),B_{\operatorname{fvar}i};
      \operatorname{VectorInputDomain},
      \operatorname{Unique}(\operatorname{fvar}i;r,b,f)).        \tag{T13}
\]
The bound-variable leaf is more substantial.  Let
\[
 u=w_i,\quad x=a_i,\quad
 E=\operatorname{entry}(b,f,w,a,i),\quad
 C=\operatorname{comparisonCode}(b,b,f,f,u),\quad D=D(r,b,f).
\]
It is exactly the following composition of requested-label providers:
\[
\begin{aligned}
 L&=\operatorname{Comp}(A,D,C;
          \operatorname{VectorInputDomain},
          \operatorname{Unique}(u;r,b,f)),\\
 K&=\operatorname{Pair}(A,E,C;
          \operatorname{VectorInputEntry}(i),L),\\
 X&=\operatorname{SubstitutionVariable}
       (\operatorname{valuePredicate}(b,f,u),x),\\
 \operatorname{Vec}(\operatorname{bvar}i)
   &=\operatorname{Comp}(A,E\mathbin{\&}C,B_{\operatorname{bvar}i};K,X).
\end{aligned}                                                   \tag{T14}
\]
For example, its uniqueness query has the precise label tuple
\[
 P_{c(A,D,C)}\bigl(P_{p(A,E,C)}(
     P_{c(A,E\mathbin{\&}C,B_{\operatorname{bvar}i})}(\lambda))\bigr).
\]
The query rank $u=w_i$ may exceed the original variable code.  It is
a call to the already completed uniqueness stratum, not a recursive
call to vector realization.  This resolves the apparent rank problem
without silently assuming a false inequality.

The entry equation required in the last branch of (T14) is exactly
(T5): substituting $x$ into the value predicate of $u$ gives the
entry formula $E$.  Together with the diagonal comparison code, this
turns the fixed substitution-variable template into
$[E\mathbin{\&}C\Rightarrow B_{\operatorname{bvar}i}]$.
Its lower closed assumption is isomorphism right transport.  Thus
(T14) does not assume the desired vector-realization theorem in a
different notation.

For the zero leaf, $B_0$ is the fixed zero comparison; pad by
$\max(B_0,[A\Rightarrow B_0])$, take its fixed proof, weaken by
$A$, and read back.  This is the same zero construction as in
coherence.  All vector leaves have now been supplied, so the
finite-demand construction on the original source term proves the
requested vector-realization theorem.

% Explicit auxiliary suppliers used by arithmetic formula and matrix closure.
\subsection{Guard collection, binder input, and canonical input}
\label{sec:term-guard-auxiliaries}

The formula and induction-matrix constructions also need three auxiliary
guard suppliers.  They do not follow merely from the arbitrary-term
vector-realization theorem, so we give their algorithms explicitly.
All translations in a native guard loop use one fixed parameter tuple
$\lambda$; no recursive call changes that tuple.

Let $G_0=\top$ and $G_{j+1}=G_j\mathbin\&E_j$ be any of the
domain or realization guards.  Given an \emph{explicit} primitive
recursive array $e_j$ of proofs of $A\to T_\lambda(E_j)$, define
$c_0$ to be truth weakening and form $c_{j+1}$ by the fixed native
implication-pairing recipe from $c_j,e_j$.  The bounded invariant is
that $c_j$ is a checked proof of $A\to T_\lambda(G_j)$, conjoined
with validity of the first $j$ entries.  Quantifier-free induction on
$j$ proves the assertion.  This is the direct-list interpretation of
\file{NumericPAeqValueDomainGuardCollection} and
\file{NumericPAeqValueRealizationGuardCollection}; it does not
extract an unspecified choice function from
$\forall j<k\,\exists e_j$.

For domain prepend, use antecedent
$T_\lambda(G_k(b))\wedge T_\lambda(D(x))$ and collect $k+1$
entries.  The zeroth proof is its right projection.  Entry $j+1$
is its left projection followed by the already completed projection
$G_k(b)\to D(b_j)$.  The list equations
$(x::b)_0=x$ and $(x::b)_{j+1}=b_j$ prove the desired conclusion
$G_{k+1}(x::b)$ exactly.  This construction is
\file{NumericPAeqValueDomainGuardPrepend.guard_prepend_certificate}.
Applying it to the bound guard and retaining the free guard gives
the input-prepend supplier used by universal matrix closure.

Write $D(x)=\operatorname{domainCode}(x)$ and
$V_x=\operatorname{variableValue}(x,x)$.
The requested self-variable proof is
$T_\lambda([D(x)\Rightarrow V_x])$.
Its fixed relational ND template proves
$D(x)\to(D(x)\wedge\operatorname{Eq}(x,x))$ from the
closed relational reflexivity sentence.  Thus it is precisely an
instance of the lower singleton-requested-template constructor,
with its sole lower nonlogical leaf supplied by the closed
\file{reflexivityNativeProof}; it needs no term recursion.
This is the exact boundary in
\file{NumericPAeqValueSelfVariableCertificate.lean}.

For binder input, let
$G=\operatorname{vectorInput}(s,r,b,f,a,w)$ and
\[
 A=G\mathbin\&D(x),\qquad
 B=\operatorname{vectorInput}
       (s+1,r+1,x::b,f,x::a,\operatorname{qVec}(w)),
 \quad S=[D(x)\Rightarrow V_x].
\]
The exact requested-label map in
\file{NumericPAeqSupportedVectorBinderInputCertificate.lean} is
\[
 d=\max(S,[A\Rightarrow B]),\qquad
 \lambda'=P_d(\lambda).
 \tag{T15}
\]
The only non-native request is the self-variable proof at $\lambda'$.
At this same tuple the new realization guard has entry zero $V_x$
and entry $j+1$ equal to the old entry $j$.  These are
\file{vector_qVec_entry_zero} and
\file{vector_qVec_entry_succ}; the successor equation is
the completed source-lift/value-syntax substitution identity, with
the exact support conditions of the vector environment.
Collect the zeroth entry using right projection and the self proof;
collect each later entry using left projection and old-guard
projection.  Prepend the domain guards of $b$ and $a$ by the
preceding native loop, retain the two copies of the free-domain
guard, and pair the resulting five factors in the literal
\file{vectorInput} association.  Native/source implication
normalization and padding readback give $T_\lambda([A\Rightarrow B])$.
All entry indices are bounded by $s+1$ or $r+1$ and every proof
value is computed explicitly.  When formula realization first
shifts $b,f,a$, this same construction is applied to those shifted
vectors; the support and validity tests are preserved by the
already completed target-shift syntax equations.

For canonical input, let $|B|=m$ and set
\begin{align*}
 G&=G_m(B),\qquad
 S_i=[D(B_i)\Rightarrow\operatorname{variableValue}(B_i,B_i)],\\
 J&=\operatorname{vectorInput}(m,0,0,B,B,\operatorname{fvarVec}(m)).
\end{align*}
Replace the source's collection-based existential bound by the
explicit primitive recursion
\[
 d_0=0,\qquad d_{i+1}=\max(d_i,S_i),\qquad
 d=d_m,\qquad \lambda'=P_d(\lambda).
 \tag{T16}
\]
The bounded invariant $\bigwedge_{i<j}S_i\le d_j$ supplies every
coverage check, including the empty case $m=0$.
At $\lambda'$ request the self-variable proof for each $i<m$;
compose it with the native projection $G\to D(B_i)$.
The canonical realization entry is literally
$\operatorname{variableValue}(B_i,B_i)$, because
$\operatorname{fvarVec}(m)_i=\operatorname{fvar}(i)$ and the
value-syntax free-variable equation selects $B_i$.
Collect these $m$ explicit arrows with the first loop.
The literal input shape is
\[
 T_{\lambda'}(J)=
 (\top\wedge T_{\lambda'}G)\wedge
   ((T_{\lambda'}G\wedge T_{\lambda'}G)\wedge
       T_{\lambda'}G_{\rm realized}).
\]
Truth weakening, three identities on $G$, and the collected arrow
therefore prove $G\to J$ by fixed conjunction operations.
Native/source normalization and padding readback finish the
requested proof.  This is the explicit realization of
\file{NumericPAeqCanonicalVectorInputCertificate.lean} and its
\file{self_codes_bounded} auxiliary.  It neither invokes an
unbounded collection principle nor recurses through an arbitrary
formula-realization supplier.

Consequently these auxiliaries belong immediately after guard
projection and the lower fixed self-variable template, before
arithmetic formula realization and induction-matrix closure.
Their only new closed leaf is ordinary relational reflexivity,
already present among the six fixed ordinal leaves of the generic
relational construction.  No edge returns to the arithmetic
formula or proof strata.

\subsection{Strata and exact lower interfaces}

The construction uses the following acyclic order:
\begin{enumerate}
\item Primitive recursive syntax, source substitution, value syntax
      (T2), and its validity/naturality joins (T3)--(T5).
\item Direct $\mathcal C_4$ logical proof constructors, requested-label
      padding/readback, and fixed relational-template instances.
\item Domain, isomorphism, and realization guard projections; the
      diagonal and repeated-domain constructions (T7)--(T8).
\item Guarded totality and guarded coherence, each by its own
      decreasing-source-code finite-demand recursion.
\item Guarded uniqueness (T11)--(T12).
\item Vector realization, with its decreasing original-source-code
      recursion and lower uniqueness calls (T13)--(T14).
\end{enumerate}
No edge returns to a higher stratum.  In particular, the fact that a
substituted or replacement term can be much larger than the original
one does not invalidate the recursion: its syntax or coherence call
belongs to a lower, already total primitive recursive stratum.

For reproducibility, the precise parameterized lower interfaces are:
\begin{itemize}
\item Native implication composition, pairing, identity, conjunction
      projections, truth weakening, and the two source/native
      implication conversions; source composition/pairing and padded
      readback use exactly (T1).
\item Variable totality at an arbitrary valid target term, duplicate
      domain for one existential payload, duplicate domain for two
      payloads, successor totality, and addition/multiplication totality
      at arbitrary valid payload formulas.
\item Variable coherence at two arbitrary valid target terms,
      zero comparison, successor coherence at two payload formulas,
      and addition/multiplication coherence at four payload formulas.
\item Substitution-variable coherence at one payload and one valid
      target term, and isomorphism reflexivity at one valid target term.
\item The fixed closed relational proof supplies for inhabited zero,
      zero-isomorphism compatibility, isomorphism reflexivity/transport/right transport, successor
      totality and compatibility, and addition/multiplication totality
      and coherence, as used by the corresponding requested templates.
\end{itemize}
These are fixed formula templates with numerical term/formula holes,
not arbitrary arithmetic-proof or PA-calibration oracles.  Their
implementation and correctness are the responsibility of the lower
generic relational stratum; this section identifies every such call
and proves that no further term supplier remains behind it.

The resulting primitive recursive functions satisfy, in
$\mathsf{iPRA}$, the desired guarded totality, guarded coherence,
uniqueness, and vector-realization implications for all inputs passing
their explicit support, environment, and requested-coverage tests.
The proof is uniform in those inputs.  It does not infer a weak-base
theorem merely from the existence of a full-PA proof in the repository.
Nor does it by itself complete any still-unproved lower relational
interface or the higher PA-proof transport.  It closes the term
supplier portion once the explicitly enumerated lower interfaces are
provided.
